% GENERATED FILE. Edit canonical lessons, then run npm run generate:books.
% canonical-source-hash: fnv1a64:b2fe2215b6a4f23a
% canonical-lessons: RU-C136-nikto, RU-C136-nichego, RU-C136-nigde, RU-C136-nikuda, RU-C136-zachem

\chapter{Nobody, Nothing, Nowhere, To Nowhere, What For}
\label{ch:ru-nobody-nothing-nowhere-to-nowhere-what-for}
\hlchaptermodality{136}

\begin{chapteropening}
\textbf{By the end of this chapter:} \emph{I can say that there is nobody and nothing, and ask what something is for.}

Complete the last lesson of chapter 136: I can say that there is nobody and nothing, and ask what something is for.
\end{chapteropening}

\section[\texorpdfstring{\ru{никто} (niktó)}{никто (niktó)}]{\ru{никто} (niktó) — nobody}

\label{lesson:RU-C136-nikto}



\begin{quote}
Before the new one: say the Russian for to draw, then the Russian for to travel.
\end{quote}

\subsection*{You'll want to know: \ru{никто}}
\textbf{\ru{никто}} (\emph{niktó}) — \textquotedblleft{}nobody\textquotedblright{}. Russian still puts \textbf{\ru{не}} before the verb: \textbf{\ru{Никто} \ru{не} \ru{знает}.} — \emph{niktó ne znáyet} — \textquotedblleft{}Nobody knows.\textquotedblright{} Two negatives make one \textquotedblleft{}no\textquotedblright{}.

\begin{grammarlens}[title={What you've built}]
One more word for this chapter.
\end{grammarlens}

\subsection*{Your turn}
\begin{itemize}
\raggedright
  \item \emph{Say it:} \emph{niktó}
  \item \emph{Say it:} \emph{niktó}, once more
  \item \emph{Say it:} say \emph{niktó}
  \item \emph{Recall:} say the Russian for to draw, then the Russian for to travel, then say \emph{niktó} again
\end{itemize}

\begin{tcolorbox}[breakable,colback=teal!4,colframe=teal!35!black,title={Before you move on}]
What does \ru{никто} mean? (\textquotedblleft{}Nobody\textquotedblright{}.) Say it once more.
\end{tcolorbox}
\section[\texorpdfstring{\ru{ничего} (nichevó)}{ничего (nichevó)}]{\ru{ничего} (nichevó) — nothing}

\label{lesson:RU-C136-nichego}



\begin{quote}
Before the new one: say the Russian for to travel, then the Russian for nobody.
\end{quote}

\subsection*{You'll want to know: \ru{ничего}}
\textbf{\ru{ничего}} (\emph{nichevó}) — \textquotedblleft{}nothing\textquotedblright{}. Spelled with \textbf{\ru{г}}, said with a \emph{v}, like \textbf{\ru{сегодня}}. \textbf{\ru{Я} \ru{ничего} \ru{не} \ru{знаю}.} — \emph{ya nichevó ne znáyu} — \textquotedblleft{}I don't know anything.\textquotedblright{} On its own, \textbf{\ru{Ничего}!} means \textquotedblleft{}It's all right!\textquotedblright{}, for instance after \textbf{\ru{Извините}!}

\begin{grammarlens}[title={What you've built}]
One more word for this chapter.
\end{grammarlens}

\subsection*{Your turn}
\begin{itemize}
\raggedright
  \item \emph{Say it:} \emph{nichevó}
  \item \emph{Say it:} \emph{nichevó}, once more
  \item \emph{Say it:} say \emph{nichevó}
  \item \emph{Recall:} say the Russian for to travel, then the Russian for nobody, then say \emph{nichevó} again
\end{itemize}

\begin{tcolorbox}[breakable,colback=teal!4,colframe=teal!35!black,title={Before you move on}]
What does \ru{ничего} mean? (\textquotedblleft{}Nothing\textquotedblright{}.) Say it once more.
\end{tcolorbox}
\section[\texorpdfstring{\ru{нигде} (nigdé)}{нигде (nigdé)}]{\ru{нигде} (nigdé) — nowhere}

\label{lesson:RU-C136-nigde}



\begin{quote}
Before the new one: say the Russian for nobody, then the Russian for nothing.
\end{quote}

\subsection*{You'll want to know: \ru{нигде}}
\textbf{\ru{нигде}} (\emph{nigdé}) — \textquotedblleft{}nowhere\textquotedblright{}. It is \textbf{\ru{где}} with \textbf{\ru{ни}-}: \textbf{\ru{Нигде} \ru{нет} \ru{такси}.} — \emph{nigdé nyet taksí} — \textquotedblleft{}There are no taxis anywhere.\textquotedblright{}

\begin{grammarlens}[title={What you've built}]
One more word for this chapter.
\end{grammarlens}

\subsection*{Your turn}
\begin{itemize}
\raggedright
  \item \emph{Say it:} \emph{nigdé}
  \item \emph{Say it:} \emph{nigdé}, once more
  \item \emph{Say it:} say \emph{nigdé}
  \item \emph{Recall:} say the Russian for nobody, then the Russian for nothing, then say \emph{nigdé} again
\end{itemize}

\begin{tcolorbox}[breakable,colback=teal!4,colframe=teal!35!black,title={Before you move on}]
What does \ru{нигде} mean? (\textquotedblleft{}Nowhere\textquotedblright{}.) Say it once more.
\end{tcolorbox}
\section[\texorpdfstring{\ru{никуда} (nikudá)}{никуда (nikudá)}]{\ru{никуда} (nikudá) — (to) nowhere}

\label{lesson:RU-C136-nikuda}



\begin{quote}
Before the new one: say the Russian for nothing, then the Russian for nowhere.
\end{quote}

\subsection*{You'll want to know: \ru{никуда}}
\textbf{\ru{никуда}} (\emph{nikudá}) — \textquotedblleft{}(to) nowhere\textquotedblright{}. It is \textbf{\ru{куда}} with \textbf{\ru{ни}-}, so it answers \textquotedblleft{}where to?\textquotedblright{}: \textbf{\ru{Я} \ru{никуда} \ru{не} \ru{иду}.} — \emph{ya nikudá ne idú} — \textquotedblleft{}I'm not going anywhere.\textquotedblright{}

\begin{grammarlens}[title={What you've built}]
One more word for this chapter.
\end{grammarlens}

\subsection*{Your turn}
\begin{itemize}
\raggedright
  \item \emph{Say it:} \emph{nikudá}
  \item \emph{Say it:} \emph{nikudá}, once more
  \item \emph{Say it:} say \emph{nikudá}
  \item \emph{Recall:} say the Russian for nothing, then the Russian for nowhere, then say \emph{nikudá} again
\end{itemize}

\begin{tcolorbox}[breakable,colback=teal!4,colframe=teal!35!black,title={Before you move on}]
What does \ru{никуда} mean? (\textquotedblleft{}(to) nowhere\textquotedblright{}.) Say it once more.
\end{tcolorbox}
\section[\texorpdfstring{\ru{зачем} (zachém)}{зачем (zachém)}]{\ru{зачем} (zachém) — what for}

\label{lesson:RU-C136-zachem}



\begin{quote}
Before the new one: say the Russian for nowhere, then the Russian for (to) nowhere.
\end{quote}

\subsection*{You'll want to know: \ru{зачем}}
\textbf{\ru{зачем}} (\emph{zachém}) — \textquotedblleft{}what for\textquotedblright{}. It asks about the purpose; \textbf{\ru{почему}} asks about the cause. \textbf{\ru{Зачем}?} — \textquotedblleft{}What for?\textquotedblright{}

\begin{grammarlens}[title={What you've built}]
One more word for this chapter.
\end{grammarlens}

\subsection*{Your turn}
\begin{itemize}
\raggedright
  \item \emph{Say it:} \emph{zachém}
  \item \emph{Say it:} \emph{zachém}, once more
  \item \emph{Say it:} say \emph{zachém}
  \item \emph{Recall:} say the Russian for nowhere, then the Russian for (to) nowhere, then say \emph{zachém} again
\end{itemize}

\begin{tcolorbox}[breakable,colback=teal!4,colframe=teal!35!black,title={Before you move on}]
What does \ru{зачем} mean? (\textquotedblleft{}What for\textquotedblright{}.) Say it once more.
\end{tcolorbox}
